\section{General production: a sharp complexity bound}\label{prod:section}
There are $n$ fixed owners and $d$ fixed resource types. Owner $i$ reports
$b_i\in\R_+^d$; write $c=\sum_i b_i$. Fix a nonempty finite family
$V=\{w^1,\ldots,w^K\}\subseteq\R_+^d$, with a report-independent priority,
and define
\[
 f(c)=\min_{w\in V}w\cdot c,\qquad
 v_b(S)=f\left(\sum_{i\in S}b_i\right),\qquad T=f(c).
\]
A rule is eligible if its payments are nonnegative, cover every coalition,
and are componentwise nondecreasing in each owner's whole bundle, with all
other reports fixed. These conditions hold on all of $(\R_+^d)^n$.
They express truthfulness for binding underreports and payment-only utility;
there are no operating costs, overreport, false-name, population-change or
coordinated-report guarantees. A factor $C$ bounds \emph{total payment} by
$CT$; external subsidy is at most $(C-1)T$. The family $V$ remains fixed
across reports. Its size may be exponential in a compact description.

\begin{theorem}[General sharp-order upper bound]\label{prod:upper}
For every integer $j\ge2$, there exists an eligible deterministic rule with factor
\begin{equation}\label{prod:constant}
 C_j(K)=j+\frac{\log K+2\log j+1}{\alpha_j},\qquad
 \alpha_j=\log\frac{j^3}{j^2+j-1}>0.
\end{equation}
With $L=\log K$ and
$j=\max\{2,\lceil L/[\log(2+L)]^2\rceil\}$, this factor is
$(1+o(1))\log K/\log\log K$, uniformly over owner counts, dimensions,
technologies and reports. Zero-value profiles receive zero; for $K=1$ the
constructed rule is efficient.
\end{theorem}

\paragraph{The mechanism.}
For $q\in(0,1)$ let independent owner bits satisfy $\Prob(B_i=0)=q$.
Choose the first minimizer of the \emph{uncapped} objective,
\[
 W(b,B)=\operatorname*{first\,argmin}_{w\in V}
               \sum_i B_i(w\cdot b_i),\qquad
 a_i(w,b)=\min\{w\cdot b_i,T\}.
\]
Only the payment contribution $a_i$ is capped. Put
\begin{align}
 z_i^q&=\E a_i(W,b),&
 p_i^q&=\E[a_i(W,b)\mid B_i=0],\label{prod:p}\\
 m_i^q&=\E_{B_{-i}}\left[f\left(b_i+\sum_{h\ne i}B_hb_h\right)
                         -f\left(\sum_{h\ne i}B_hb_h\right)\right],
 &r_i^q&=qp_i^q+(1-q)m_i^q.\label{prod:hybrid}
\end{align}
The final payment is the deterministic finite expectation
\begin{equation}\label{prod:payment}
 M_i(b)=j\int_0^{1/j}r_i^q(b)\,dq.
\end{equation}
These quantities are finite polynomials in $q$ at a fixed profile; no
continuity or generic-position assumption about the reports is required.

\begin{proof}[Proof of Theorem~\ref{prod:upper}]
Conditional on $B_i=0$, every score and the selected form are independent of
$b_i$. For that fixed nonnegative form, both $w\cdot b_i$ and $T$ are
nondecreasing in $b_i$, so their minimum is nondecreasing. In $m_i^q$ the
off term is independent of $b_i$ and the on term is nondecreasing. The
nonnegative fixed mixtures preserve this property and nonnegativity.

For any $w,S$, if some $i\in S$ has $w\cdot b_i\ge T$, then
$\sum_{i\in S}a_i(w,b)\ge T\ge v_b(S)$. Otherwise no term in $S$ is capped,
and the sum is $w\cdot\sum_{i\in S}b_i\ge v_b(S)$. Thus every capped form
covers coalitions. Fix $B_{-i}$ and let $w_1$ be the on minimizer. The off
optimum is at most the off objective of this same form, whence
\begin{equation}\label{prod:endpoint}
 f\left(b_i+\sum_{h\ne i}B_hb_h\right)-f\left(\sum_{h\ne i}B_hb_h\right)
 \ge w_1\cdot b_i=a_i(w_1,b).
\end{equation}
The equality holds because the contribution is a nonnegative summand of
the selected on value, which is at most $T$. Averaging proves $r_i^q\ge z_i^q$,
and hence coverage for $r^q$ and $M$.

For the budget define $g(q)=\E f(\sum_iB_ib_i)$, $Z(q)=\sum_i z_i^q$,
and $R(q)=\sum_i r_i^q$. Capping leaves every open contribution unchanged:
\[
 \sum_iB_i a_i(W,b)=f\left(\sum_iB_ib_i\right).
\]
Consequently $q\sum_i p_i^q=Z-g$. Differentiating the finite mask polynomial
gives $\sum_i m_i^q=-g'$, so $R=Z-g-(1-q)g'$. Since $g(0)=T$ and $g\ge0$,
for $0<\eta<1$ integration by parts gives the exact identity
\begin{equation}\label{prod:integral}
 \int_0^\eta R\,dq=\int_0^\eta Z\,dq+T-(1-\eta)g(\eta)
                         -2\int_0^\eta g\,dq
 \le\int_0^\eta Z\,dq+T.
\end{equation}
Both the sign and the factor two in this identity matter.

Suppose $T>0$. For each fixed form let
$d_i(w)=a_i(w,b)/T\in[0,1]$ and $\rho_w=\sum_i d_i(w)$.
The selected form satisfies $\sum_iB_id_i(W)\le1$. For $\lambda>0$ put
$\alpha=-\log(q+(1-q)e^{-\lambda})$. The convex function
$u\mapsto\log(q+(1-q)e^{-\lambda u})$ lies below its endpoint chord
$-\alpha u$ on $[0,1]$. Independence and Markov's inequality give
\[
 \Prob\left(\sum_iB_id_i(w)\le1\right)
 \le e^\lambda\prod_i(q+(1-q)e^{-\lambda d_i(w)})
 \le e^{\lambda-\alpha\rho_w}.
\]
A union over the fixed family and integration of the resulting tail yield
\begin{equation}\label{prod:chernoff}
 \Prob(\rho_W\ge x)\le\min\{1,K e^{\lambda-\alpha x}\},\qquad
 \frac{Z(q)}T\le\frac{\log K+\lambda+1}{\alpha}.
\end{equation}
Take $\eta=1/j$ and $\lambda=2\log j$. For $q\le1/j$,
$q+(1-q)e^{-\lambda}\le(j^2+j-1)/j^3$, so
\eqref{prod:integral}--\eqref{prod:chernoff} give \eqref{prod:constant}.
For the displayed choice of $j$, as $L\to\infty$,
$j\sim L/(\log L)^2$, $\alpha_j\sim\log L$, and
$j=o(L/\log L)$, proving the asserted asymptotic.

If $T=0$, all masked production values and capped contributions vanish.
If $K=1$, no contribution is capped, $p_i^q=m_i^q=w\cdot b_i$, and
$M_i=w\cdot b_i$. Finally each finite-mask term is
$q^{n-1-|S|}(1-q)^{|S|}[q a_i(W(b,1_S),b)+(1-q)(v_b(S\cup\{i\})-v_b(S))]$.
Its integral is a finite sum with rational coefficients when $j$ is integer.
This gives exact rational payments for rational input, without an efficiency
claim about enumeration of masks or forms.
\end{proof}

\begin{corollary}[Production LPs and networks]\label{prod:lp}
The theorem applies to
$f(c)=\max\{\pi\cdot x:x\ge0,Ax\le c\}$ for fixed $A\ge0$, with every
activity column nonzero and fixed positive profits $\pi$. Here $V$ can be
the vertices of $\{w\ge0:A^Tw\ge\pi\}$. For single-commodity directed flow,
one may instead use the characteristic vectors of the inclusion-minimal
structural source--sink edge cuts.
\end{corollary}
\begin{proof}
Each activity is bounded by a positive entry in its column, so every feasible
production set is compact at finite $c$. The dual is nonempty, pointed and
polyhedral. Strong duality and attainment at a vertex give the claimed finite
min-linear representation, including boundary capacities. Network max-flow/
min-cut gives the cut representation directly. Redundant nonminimal cuts may
be removed once and for all because capacities are nonnegative. All owners
and forms remain fixed across profiles. A positive-profit zero-resource
activity is excluded because it would make production unbounded.
\end{proof}
Dual core allocations are established ingredients~\cite{Owen,SametZemel}.
General dual vertices may be fractional; their count is not an ordinary
network-cut count. Constraints introducing capacity-independent dual offsets
are outside the homogeneous model here.

\begin{proposition}[Why the payment cap is essential]\label{prod:uncapped}
Removing the cap from the own-excluded contribution in
\eqref{prod:p} makes the uniform-$q$ hybrid unbounded even with two owners
and two supporting forms.
\end{proposition}
\begin{proof}
Take unit-profit activity columns $(1,1)$ and $(M,0)$, $M>1$, and reports
$c=(M,1)$. The two dual vertices are $(1,0)$ and $(1/M,1-1/M)$, and
$T=2-1/M$. If owner 1 is off and owner 2 on, $(1,0)$ uniquely minimizes the
raw score, giving uncapped contribution $M$. Thus $r_1^q\ge q(1-q)M$;
with $j=2$ its average is at least $M/6$, an unbounded ratio to $T$.
This obstructs only the uncapped rule: since
$f(c)=c_1/M+(1-1/M)\min(c_1,c_2)$, paying the linear part to owner 1 and
splitting the remaining term equally gives an efficient eligible rule.
\end{proof}
